\documentclass[11pt,a4paper]{article}
\usepackage[utf8]{inputenc}
\usepackage[french]{babel}
\usepackage[T1]{fontenc}
\usepackage{amsmath,amssymb}
\usepackage{graphicx}
\usepackage{xcolor}
\usepackage{hyperref}
\usepackage[margin=2cm]{geometry}
\usepackage{tabularx}
\usepackage{array}
\usepackage{colortbl}
\usepackage{tcolorbox}

% Définition des couleurs
\definecolor{titrebleu}{RGB}{0,102,204}
\definecolor{defvert}{RGB}{0,153,76}
\definecolor{exempleorange}{RGB}{255,102,0}
\definecolor{importantrouge}{RGB}{204,0,0}
\definecolor{fondjaune}{RGB}{255,250,205}
\definecolor{fondvert}{RGB}{230,255,230}
\definecolor{fondbleu}{RGB}{230,240,255}
\definecolor{fondgris}{RGB}{245,245,245}

% Boîtes colorées
\newtcolorbox{important}{
    colback=fondjaune,
    colframe=exempleorange,
    fonttitle=\bfseries,
    title=Important,
    arc=3mm
}

\newtcolorbox{conseil}{
    colback=fondvert,
    colframe=defvert,
    fonttitle=\bfseries,
    title=Conseil,
    arc=3mm
}

\newtcolorbox{info}{
    colback=fondbleu,
    colframe=titrebleu,
    fonttitle=\bfseries,
    title=Information,
    arc=3mm
}

% Style des sections
\usepackage{titlesec}
\titleformat{\section}
{\color{titrebleu}\normalfont\LARGE\bfseries}
{\color{titrebleu}\thesection}{1em}{}

\titleformat{\subsection}
{\color{defvert}\normalfont\Large\bfseries}
{\color{defvert}\thesubsection}{1em}{}

\title{\textbf{\Huge\color{titrebleu}Guide GitHub Student Pack}\\[0.5em]
\large\color{defvert}Carte de Navigation Complète\\
Pour Étudiants en Médecine et Informatique}
\author{\large Guide Pratique}
\date{\large Année 2024-2025}

\begin{document}

\maketitle

\begin{info}
\textbf{Qu'est-ce que le GitHub Student Pack ?}

Le GitHub Student Pack est un ensemble d'outils et services gratuits offerts aux étudiants. Il comprend plus de 100 offres dans différents domaines : développement, design, déploiement, apprentissage, etc.

\textbf{Comment l'obtenir ?}
\begin{enumerate}
    \item Aller sur \url{https://education.github.com/pack}
    \item S'inscrire avec votre email étudiant (@univ-*.fr)
    \item Valider votre statut d'étudiant
\end{enumerate}
\end{info}

\tableofcontents
\newpage

\section{Organisation et Productivité}

\subsection{Notion + Notion Template Collection}

\textbf{Ce que c'est :} Espace de travail tout-en-un pour notes, bases de données, wikis, tableaux Kanban.

\textbf{Offre Student Pack :}
\begin{itemize}
    \item Plan Plus gratuit (valeur 8\$/mois)
    \item Collection de templates étudiants
\end{itemize}

\textbf{Pourquoi c'est utile :}
\begin{itemize}
    \item Centraliser cours, fiches de révision, QCM
    \item Planning de révisions et projets
    \item Organisation cours médecine + informatique au même endroit
    \item Templates prêts à l'emploi
\end{itemize}

\begin{conseil}
Idéal pour créer une base de connaissances personnelle avec système de tags et recherche rapide.
\end{conseil}

\subsection{PomoDone + HazeOver}

\textbf{PomoDone :} Application Pomodoro (technique 25 min travail + 5 min pause)

\textbf{HazeOver :} Réduit les distractions en assombrissant les fenêtres en arrière-plan

\textbf{Pourquoi c'est utile :}
\begin{itemize}
    \item Maintenir la concentration pendant longues sessions
    \item Parfait pour révisions médecine et sessions de code
    \item Éviter la fatigue mentale
\end{itemize}

\newpage
\section{Sécurité et Gestion de Comptes}

\subsection{1Password + Dashlane}

\textbf{Ce que c'est :} Gestionnaires de mots de passe sécurisés

\textbf{Offre Student Pack :}
\begin{itemize}
    \item 1Password : 6 mois gratuits
    \item Dashlane : 1 an gratuit
\end{itemize}

\textbf{Pourquoi c'est utile :}
\begin{itemize}
    \item Beaucoup d'outils = beaucoup de comptes
    \item Mots de passe forts et uniques automatiquement
    \item Gain de temps (auto-remplissage)
    \item Protection contre le piratage
\end{itemize}

\begin{important}
Avec tous les outils du Student Pack, vous aurez des dizaines de comptes. Un gestionnaire de mots de passe est ESSENTIEL.
\end{important}

\subsection{Doppler}

\textbf{Ce que c'est :} Gestion de secrets et variables d'environnement

\textbf{Pourquoi c'est utile :}
\begin{itemize}
    \item Gérer clés API, tokens, mots de passe de BDD
    \item Travailler en équipe sans partager secrets par email
    \item Synchronisation automatique entre environnements
\end{itemize}

\begin{conseil}
Plus avancé. Utile quand vous gérez des projets avec bases de données ou APIs externes.
\end{conseil}

\newpage
\section{Développement et Programmation}

\subsection{GitHub Copilot}

\textbf{Ce que c'est :} Assistant IA qui aide à écrire du code

\textbf{Offre Student Pack :} Gratuit pendant vos études

\textbf{Pourquoi c'est utile :}
\begin{itemize}
    \item Autocomplétion intelligente du code
    \item Explications de code existant
    \item Génération de tests
    \item Correction de bugs
    \item Gain de temps énorme en TD et projets
\end{itemize}

\begin{important}
L'un des outils les PLUS utiles du pack pour les cours d'informatique. À activer en priorité !
\end{important}

\subsection{JetBrains Suite}

\textbf{Ce que c'est :} Suite d'IDEs professionnels (PyCharm, IntelliJ IDEA, WebStorm, etc.)

\textbf{Offre Student Pack :} Tous les IDEs gratuits

\textbf{Pourquoi c'est utile :}
\begin{itemize}
    \item Environnements de développement professionnels
    \item Débogueur puissant
    \item Refactoring automatique
    \item PyCharm pour Python, IntelliJ pour Java
\end{itemize}

\subsection{Visual Studio Code}

\textbf{Ce que c'est :} Éditeur de code léger et extensible

\textbf{Pourquoi c'est utile :}
\begin{itemize}
    \item Plus léger que JetBrains
    \item Extensions pour tous les langages
    \item Intégration Git native
    \item Gratuit de base, mais extensions premium avec Student Pack
\end{itemize}

\subsection{GitHub Codespaces}

\textbf{Ce que c'est :} Environnement de développement dans le cloud

\textbf{Offre Student Pack :} 180 heures/mois gratuites

\textbf{Pourquoi c'est utile :}
\begin{itemize}
    \item Coder depuis n'importe quel ordinateur
    \item Pas d'installation locale nécessaire
    \item Configuration automatique de l'environnement
    \item Idéal si vous changez souvent de PC (fac/maison)
\end{itemize}

\subsection{GitHub Pro}

\textbf{Ce que c'est :} Version premium de GitHub

\textbf{Offre Student Pack :} Gratuit pendant vos études

\textbf{Fonctionnalités :}
\begin{itemize}
    \item Repos privés illimités
    \item GitHub Pages sur repos privés
    \item Actions minutes augmentées
    \item Insights avancés
\end{itemize}

\newpage
\section{Gestion de Versions (Git)}

\subsection{GitLens}

\textbf{Ce que c'est :} Extension VS Code pour visualiser l'historique Git

\textbf{Pourquoi c'est utile :}
\begin{itemize}
    \item Voir qui a modifié chaque ligne
    \item Historique des fichiers clair
    \item Navigation dans les commits
\end{itemize}

\subsection{GitHub Desktop + GitKraken + Tower}

\textbf{Ce que c'est :} Interfaces graphiques pour Git

\textbf{Pourquoi c'est utile :}
\begin{itemize}
    \item Visualiser branches et merges facilement
    \item Pas besoin de commandes en ligne
    \item Éviter les erreurs Git
    \item Idéal pour débutants
\end{itemize}

\begin{conseil}
\textbf{Recommandation :}
\begin{itemize}
    \item Débutant : GitHub Desktop (simple)
    \item Intermédiaire : GitKraken (visuel)
    \item Avancé : Tower (puissant)
\end{itemize}
\end{conseil}

\subsection{Working Copy (mobile)}

\textbf{Ce que c'est :} Client Git pour iOS

\textbf{Pourquoi c'est utile :}
\begin{itemize}
    \item Modifier code depuis téléphone/tablette
    \item Faire commits en déplacement
    \item Utile pour petites corrections urgentes
\end{itemize}

\newpage
\section{Apprentissage et Formation}

\subsection{DataCamp}

\textbf{Ce que c'est :} Plateforme d'apprentissage data science et programmation

\textbf{Offre Student Pack :} 3 mois gratuits

\textbf{Pourquoi c'est utile :}
\begin{itemize}
    \item Cours Python, R, SQL
    \item Data analysis et statistiques
    \item Machine Learning
    \item Très pertinent pour projets santé/bio-informatique
\end{itemize}

\begin{conseil}
EXCELLENT pour médecine + informatique : analyse de données médicales, statistiques, recherche.
\end{conseil}

\subsection{Educative}

\textbf{Ce que c'est :} Cours de programmation interactifs

\textbf{Offre Student Pack :} 6 mois gratuits

\textbf{Pourquoi c'est utile :}
\begin{itemize}
    \item Cours structurés avec exercices
    \item Pas besoin d'installer quoi que ce soit
    \item Sujets variés : web, algorithmes, systèmes
\end{itemize}

\subsection{Codédex}

\textbf{Ce que c'est :} Apprentissage ludique de la programmation

\textbf{Pourquoi c'est utile :}
\begin{itemize}
    \item Gamification de l'apprentissage
    \item Motivation par badges et niveaux
    \item Bon pour débutants
\end{itemize}

\subsection{Scrimba + FrontendMasters}

\textbf{Ce que c'est :} Formations développement web

\textbf{Scrimba :} Vidéos interactives où on peut modifier le code en direct

\textbf{FrontendMasters :} Cours avancés front-end et JavaScript

\textbf{Pourquoi c'est utile :}
\begin{itemize}
    \item Créer des sites web modernes
    \item React, Vue, Angular
    \item Portfolio professionnel
\end{itemize}

\newpage
\section{Entraînement Algorithmes}

\subsection{AlgoExpert + Interview Cake}

\textbf{Ce que c'est :} Plateformes d'entraînement aux algorithmes et entretiens techniques

\textbf{Offre Student Pack :}
\begin{itemize}
    \item AlgoExpert : réduction
    \item Interview Cake : accès premium
\end{itemize}

\textbf{Pourquoi c'est utile :}
\begin{itemize}
    \item Préparer entretiens techniques
    \item Stages et alternances en informatique
    \item Améliorer logique algorithmique
\end{itemize}

\begin{conseil}
Prioritaire si vous visez stages/alternance en entreprise tech. Moins urgent si objectif = seulement réussir examens de fac.
\end{conseil}

\newpage
\section{Data Science et Intelligence Artificielle}

\subsection{Deepnote}

\textbf{Ce que c'est :} Notebooks collaboratifs (comme Jupyter) dans le cloud

\textbf{Offre Student Pack :} Plan éducation gratuit

\textbf{Pourquoi c'est utile :}
\begin{itemize}
    \item Analyse de données
    \item Visualisations interactives
    \item Collaboration en temps réel
    \item Machine Learning
\end{itemize}

\begin{conseil}
TRÈS pertinent pour projets santé : analyse données patients, stats biomédicales, recherche.
\end{conseil}

\subsection{Microsoft Azure}

\textbf{Ce que c'est :} Cloud computing (serveurs, IA, bases de données)

\textbf{Offre Student Pack :} 100\$ de crédits

\textbf{Pourquoi c'est utile :}
\begin{itemize}
    \item Services d'IA (vision, NLP)
    \item Hébergement applications
    \item Calcul haute performance
    \item Bases de données
\end{itemize}

\newpage
\section{Hébergement et Déploiement}

\subsection{GitHub Pages}

\textbf{Ce que c'est :} Hébergement gratuit de sites statiques

\textbf{Pourquoi c'est utile :}
\begin{itemize}
    \item Portfolio personnel
    \item Documentation de projets
    \item Sites vitrines
    \item Gratuit et illimité
\end{itemize}

\subsection{Noms de Domaine (Name.com, Namecheap, .TECH)}

\textbf{Ce que c'est :} Votre propre nom de domaine

\textbf{Offre Student Pack :}
\begin{itemize}
    \item Name.com : 1 domaine .me gratuit
    \item Namecheap : 1 an gratuit .me
    \item .TECH : 1 domaine .tech gratuit
\end{itemize}

\textbf{Pourquoi c'est utile :}
\begin{itemize}
    \item Portfolio professionnel : prenom-nom.tech
    \item Plus crédible qu'un sous-domaine GitHub
    \item Email personnalisé possible
\end{itemize}

\subsection{Heroku + DigitalOcean}

\textbf{Ce que c'est :} Hébergement d'applications avec backend

\textbf{Offre Student Pack :}
\begin{itemize}
    \item Heroku : crédits gratuits
    \item DigitalOcean : 200\$ de crédits
\end{itemize}

\textbf{Pourquoi c'est utile :}
\begin{itemize}
    \item Déployer applications Django, Flask, Node.js
    \item APIs
    \item Applications avec base de données
\end{itemize}

\subsection{MongoDB Atlas + Appwrite}

\textbf{Ce que c'est :} Bases de données dans le cloud

\textbf{MongoDB Atlas :} Base NoSQL

\textbf{Appwrite :} Backend as a Service (BDD + Auth + Storage)

\textbf{Pourquoi c'est utile :}
\begin{itemize}
    \item Pas besoin de gérer serveur
    \item Scalabilité automatique
    \item Backend prêt rapidement
\end{itemize}

\newpage
\section{Tests et Qualité}

\subsection{BrowserStack + LambdaTest + Polypane}

\textbf{Ce que c'est :} Tests multi-navigateurs et appareils

\textbf{Pourquoi c'est utile :}
\begin{itemize}
    \item Tester site sur Safari/iOS sans iPhone
    \item Compatibilité tous navigateurs
    \item Tests responsive design
\end{itemize}

\begin{conseil}
Utile quand votre site devient public et utilisé par d'autres personnes.
\end{conseil}

\subsection{Travis CI + Codecov + DeepScan}

\textbf{Travis CI :} Intégration continue (tests automatiques)

\textbf{Codecov :} Couverture de tests

\textbf{DeepScan :} Analyse qualité JavaScript

\textbf{Pourquoi c'est utile :}
\begin{itemize}
    \item Tests automatiques à chaque push
    \item Détecter bugs avant mise en production
    \item Qualité de code
\end{itemize}

\subsection{Sentry + Honeybadger + New Relic + Datadog}

\textbf{Ce que c'est :} Monitoring d'applications en production

\textbf{Pourquoi c'est utile :}
\begin{itemize}
    \item Détecter erreurs en temps réel
    \item Monitoring des performances
    \item Logs centralisés
\end{itemize}

\begin{conseil}
Niveau production/professionnel. Utile quand vous avez des utilisateurs réels. Pour TD : souvent inutile.
\end{conseil}

\newpage
\section{Design et Création}

\subsection{Icons8 + IconScout + Octicons}

\textbf{Ce que c'est :} Bibliothèques d'icônes

\textbf{Offre Student Pack :} Accès premium

\textbf{Pourquoi c'est utile :}
\begin{itemize}
    \item Rendre projets plus professionnels
    \item Interfaces utilisateur
    \item Présentations
\end{itemize}

\subsection{Bootstrap Studio}

\textbf{Ce que c'est :} Outil visuel pour créer sites Bootstrap

\textbf{Pourquoi c'est utile :}
\begin{itemize}
    \item Créer sites rapidement
    \item Design responsive automatique
    \item Pas besoin d'écrire tout le HTML/CSS
\end{itemize}

\subsection{Visme}

\textbf{Ce que c'est :} Création de présentations, infographies, rapports

\textbf{Pourquoi c'est utile :}
\begin{itemize}
    \item Présentations professionnelles
    \item Exposés de projets
    \item Infographies pour rapports
\end{itemize}

\subsection{ToDiagram}

\textbf{Ce que c'est :} Création de diagrammes et schémas

\textbf{Pourquoi c'est utile :}
\begin{itemize}
    \item Schématiser algorithmes
    \item Architecture d'applications
    \item Diagrammes de révision (médecine)
    \item Organigrammes
\end{itemize}

\newpage
\section{Outils Spécialisés}

\subsection{POEditor}

\textbf{Ce que c'est :} Gestion de traductions pour applications

\textbf{Pourquoi c'est utile :}
\begin{itemize}
    \item Applications multilingues
    \item Collaboration avec traducteurs
\end{itemize}

\begin{conseil}
Non prioritaire sauf si vous créez une app multilingue.
\end{conseil}

\subsection{Stripe}

\textbf{Ce que c'est :} Gestion de paiements en ligne

\textbf{Offre Student Pack :} Frais réduits

\textbf{Pourquoi c'est utile :}
\begin{itemize}
    \item Projets commerciaux
    \item Abonnements
    \item Marketplace
\end{itemize}

\begin{conseil}
Uniquement si projet commercial. Sinon pas nécessaire.
\end{conseil}

\subsection{Zyte + Requestly}

\textbf{Zyte :} Web scraping (collecte de données web)

\textbf{Requestly :} Debug HTTP, modification de requêtes

\textbf{Pourquoi c'est utile :}
\begin{itemize}
    \item Collecte automatique de données
    \item Analyse de sites web
    \item Debug avancé
\end{itemize}

\begin{important}
Attention aux règles légales et CGU des sites lors du scraping !
\end{important}

\subsection{Arduino + Adafruit}

\textbf{Ce que c'est :} Électronique et IoT

\textbf{Offre Student Pack :} Réductions sur matériel

\textbf{Pourquoi c'est utile :}
\begin{itemize}
    \item Projets hardware
    \item Systèmes embarqués
    \item Robotique
\end{itemize}

\newpage
\section{Communauté et Certifications}

\subsection{GitHub Community Exchange}

\textbf{Ce que c'est :} Plateforme pour trouver des projets open source étudiants

\textbf{Pourquoi c'est utile :}
\begin{itemize}
    \item Collaborer avec autres étudiants
    \item Contribuer à projets réels
    \item Construire portfolio
    \item Networking
\end{itemize}

\subsection{GitHub Campus Experts}

\textbf{Ce que c'est :} Programme de leadership technique

\textbf{Pourquoi c'est utile :}
\begin{itemize}
    \item Former communauté tech dans votre fac
    \item Organiser événements
    \item Accès ressources exclusives
    \item Reconnaissance officielle
\end{itemize}

\subsection{GitHub Certification Offer 2025}

\textbf{Ce que c'est :} Certifications officielles GitHub

\textbf{Offre Student Pack :} Examens gratuits ou réduits

\textbf{Pourquoi c'est utile :}
\begin{itemize}
    \item Valider compétences Git/GitHub
    \item Badge officiel sur CV
    \item Crédibilité professionnelle
\end{itemize}

\newpage
\section{Guide de Priorités}

\begin{table}[h!]
\centering
\renewcommand{\arraystretch}{1.3}
\begin{tabularx}{\textwidth}{|l|X|}
\hline
\rowcolor{fondjaune}
\textbf{PRIORITÉ HAUTE} & \textbf{À activer dès maintenant} \\
\hline
GitHub Copilot & Aide énorme pour coder \\
GitHub Pro & Fonctionnalités essentielles \\
1Password/Dashlane & Sécurité des comptes \\
JetBrains/VS Code & Environnement de dev \\
GitHub Desktop & Interface Git simple \\
DataCamp & Apprentissage data/Python \\
\hline
\rowcolor{fondvert}
\textbf{PRIORITÉ MOYENNE} & \textbf{Activer selon besoins} \\
\hline
Notion & Si besoin organisation \\
GitHub Pages + Domaine & Si vous créez portfolio \\
Deepnote & Si projets data \\
Educative & Si apprentissage structuré \\
Icons8/Visme & Si présentations \\
\hline
\rowcolor{fondbleu}
\textbf{PRIORITÉ BASSE} & \textbf{Activer plus tard} \\
\hline
AlgoExpert & Seulement si entretiens tech \\
Heroku/DigitalOcean & Quand app à déployer \\
BrowserStack & Quand site public \\
Sentry/Monitoring & Niveau production \\
Arduino/Adafruit & Si projets hardware \\
\hline
\end{tabularx}
\end{table}

\newpage
\section{Récapitulatif par Profil}

\subsection{Profil : Débutant en Programmation}

\textbf{Outils recommandés :}
\begin{enumerate}
    \item GitHub Copilot (aide à écrire)
    \item GitHub Desktop (Git facile)
    \item VS Code (éditeur simple)
    \item DataCamp ou Codédex (apprentissage)
    \item 1Password (sécurité)
\end{enumerate}

\subsection{Profil : Étudiant Médecine + Informatique}

\textbf{Outils recommandés :}
\begin{enumerate}
    \item Notion (organisation cours)
    \item DataCamp (data science santé)
    \item Deepnote (analyse données)
    \item GitHub Copilot (projets)
    \item Visme (présentations)
\end{enumerate}

\subsection{Profil : Projets Web}

\textbf{Outils recommandés :}
\begin{enumerate}
    \item GitHub Copilot
    \item VS Code + Extensions
    \item GitHub Pages + Domaine
    \item Bootstrap Studio
    \item Icons8/IconScout
    \item Heroku (déploiement)
\end{enumerate}

\subsection{Profil : Data Science / IA}

\textbf{Outils recommandés :}
\begin{enumerate}
    \item DataCamp
    \item Deepnote
    \item GitHub Copilot
    \item Azure (crédits IA)
    \item MongoDB Atlas
\end{enumerate}

\newpage
\section{Conseils Finaux}

\begin{important}
\textbf{N'activez pas tout d'un coup !}

Stratégie recommandée :
\begin{enumerate}
    \item Semaine 1 : Activer les 5 outils prioritaires
    \item Mois 1 : Apprendre à les utiliser correctement
    \item Mois 2+ : Ajouter outils selon besoins réels
\end{enumerate}
\end{important}

\begin{conseil}
\textbf{Astuces :}
\begin{itemize}
    \item Créez un document pour noter vos identifiants
    \item Utilisez 1Password/Dashlane dès le début
    \item Lisez la documentation de chaque outil activé
    \item Commencez par des projets simples
    \item N'hésitez pas à demander de l'aide sur les forums
\end{itemize}
\end{conseil}

\subsection{Ressources Supplémentaires}

\textbf{Documentation officielle :}
\begin{itemize}
    \item GitHub Education : \url{https://education.github.com}
    \item Student Pack : \url{https://education.github.com/pack}
    \item GitHub Docs : \url{https://docs.github.com}
\end{itemize}

\textbf{Communauté :}
\begin{itemize}
    \item GitHub Community Forum
    \item Discord GitHub Education
    \item Stack Overflow
\end{itemize}

\vspace{2cm}

\begin{center}
\textcolor{titrebleu}{\Large\textbf{Fin du Guide}}

\vspace{0.5cm}

\textcolor{defvert}{\large Bon développement et bon apprentissage !}

\vspace{1cm}

\textcolor{importantrouge}{\small N'oubliez pas : Le Student Pack est valide pendant toutes vos études.\\
Profitez-en au maximum !}
\end{center}

\end{document}